\medskip

Pour chacune des affirmations suivantes, indiquer si elle est vraie ou fausse. Chaque
réponse doit être justifiée. Une réponse non justifiée ne rapporte aucun point.

\medskip

	\begin{enumerate}
		\item On considère ci-dessous le tableau de variations d'une fonction $f$ définie sur $\R\ \{- 2\}$.

\begin{center}
\psset{xunit=1.5cm,arrowsize=2pt 3}
\begin{pspicture}(7,2.5)
\psframe(7,2.5)\psline(0,2)(7,2)\psline(1,0)(1,2.5)
\psline(2.95,0)(2.95,2)\psline(3.05,0)(3.05,2)
\uput[u](0.5,1.9){$x$} \uput[u](1.4,1.9){$- \infty$} \uput[u](3,1.9){$- 2$} \uput[u](5,1.9){$1$} \uput[u](6.6,1.9){$+ \infty$}
\psline{->}(1.5,1.5)(2.5,0.5)\psline{->}(3.5,0.5)(4.5,1.5)\psline{->}(5.5,1.5)(6.5,0.5)
\uput[d](1.2,2){5}\uput[u](2.6,0){$-\infty$}\uput[u](3.4,0){$- \infty$}\uput[d](5,2){3}\uput[u](6.8,0){1}
\rput(0.5,1){$f$}
\end{pspicture}
\end{center}



	\begin{enumerate}
		\item \textbf{Affirmation 1 :}

La droite d'équation $y = -2$ est asymptote horizontale à la courbe $\mathcal{C}_f$ de la fonction $f$.
		\item \textbf{Affirmation 2:}

$\displaystyle\lim_{x \to - \infty} \dfrac{2}{f(x) - 5} = + \infty$.
	\end{enumerate}
\item On considère la fonction $g$ définie sur $\R$ par $g(x) = x\e^{-x}$.
	\begin{enumerate}
		\item \textbf{Affirmation 3 :}
		
Le point A$\left(2~;~\dfrac{2}{\e^2}\right)$ est l'unique point d'inflexion de la courbe $\mathcal{C}_g$ de la fonction~$g$.
		\item \textbf{Affirmation 4 :}
		
Pour tout nombre réel $x$ appartenant à $] - \infty~;~2[$, on a $g(x) \leqslant x$.
	\end{enumerate}
\item \textbf{Affirmation 5 :}

L'équation $x \ln (x) = 1$ admet exactement deux solutions sur l'intervalle $]0~;~ +\infty[$.
\end{enumerate}
